\subsection{What a finite quotient can, and cannot, test}

\begin{claim}{B.f1}{Refuted}{An arbitrary two-generator finite quotient does not carry the full system}
A surjection $F_2\to G$ alone does not define the substitutions
$f(x,y)\mapsto f(y,x)$ and $f(x,y)\mapsto f(y,z)$ on elements of $G$.
Nor does it define a pentagon test. A correct finite test either retains
the word and all its evaluations, or uses a common invariant quotient for
all required maps. In particular the Legendre images at even level $n$
are not the full groups $\SL_2(\mathbb Z/n)$ requested literally in B(e).
\end{claim}
% anchor:B.f1
\manuscriptanchor{\S \S 38--39; fols.297--300; scans18--24}

\begin{proof}
In $S_3$ choose $x=(12)$ and $y=(123)$. The relation $x^2=1$ holds,
but its image under swapping the two generators is $y^2\ne1$.
Thus the swapping map does not descend. This counterexample is checked
in \nolinkurl{finite_congruence.py}. The pentagon has five evaluations in
a quotient of the geometric group of $M_{0,5}$; one quotient of $F_2$
does not specify that target or its five structure maps. Finally both
Legendre matrices $L_0=\left(\begin{smallmatrix}1&-2\\0&1\end{smallmatrix}\right)$,
$L_1=\left(\begin{smallmatrix}1&0\\2&1\end{smallmatrix}\right)$ are
identity modulo $2$, whereas $\SL_2(\mathbb F_2)$ has order six.
The projective image is the kernel of reduction to $S_3$ at even level,
as follows from the index-six modular covering and reduction; its orders
are computed below. These are failures of the proposed formulation, not
counterexamples to a correctly formulated Grothendieck--Teichm\"uller relation.
\end{proof}

For a finite two-generated $G$, put
\[
 N_G=\bigcap_{q:F_2\twoheadrightarrow G}\ker q,
 \qquad \overline G=F_2/N_G.
\]
The intersection is finite. It is invariant under all automorphisms of
$F_2$, including the swapping and cyclic substitutions. More strongly,
any generating pair of $\overline G$ is the image of its distinguished
pair under an automorphism: every epimorphism to $G$ is accounted for in
the intersection. We use the finite shadow $\mathcal D(G)$ consisting of
outer automorphisms of $\overline G$ that send $x$ to a conjugate of
$x^k$, with $k$ prime to its order, and commute in the outer group with
$\theta:(x,y)\mapsto(y,x)$ and
$\delta:(x,y)\mapsto((xy)^{-1},y)$. This is Guillot's $GT(G)$, whose
inverse limit is $\GTz$, not his notation for the pentagon group
\citep[\S\S2.1--2.4]{Guillot2015root}. Its subgroup with a representative
$k=1$ is denoted $\mathcal D_1(G)$. At a finite level the multiplier need
not be unique, so this subgroup is not automatically the image of the
global subgroup $\lambda=1$.

\begin{claim}{B.f2}{Proved}{Complete matrix residue enumeration at the requested levels}
Let $Q_n=\SL_2(\mathbb Z/n)/\langle-I\rangle$ and
$H_n=\langle L_0,L_1\rangle\subseteq Q_n$. Write
$R=\left(\begin{smallmatrix}0&1\\-1&1\end{smallmatrix}\right)$ and
$S=\left(\begin{smallmatrix}0&1\\-1&0\end{smallmatrix}\right)$, so that
$\rho=\Ad R$, $\sigma=\Ad S$ on $H_n$. The following is an exhaustive
enumeration of
\[
 \xi\rho(\xi)\rho^2(\xi)=1,\quad
 \eta\sigma(\eta)=1,\quad \xi=\eta L_1^\nu,
 \qquad 2\nu+1\in(\mathbb Z/n)^\times,
 \tag{B.f2.1}
\]
with $\nu$ taken modulo the order of $L_1$. The columns headed cycles
enumerate (I),(II) with $f\in H_n'$ in the convention of B.m2.
\end{claim}
% anchor:B.f2
\manuscriptanchor{\S 46; fols.393--394,407--408; scans189--190,215--217}

\begin{center}\small
\begin{tabular}{rrrrrrrr}
\toprule
$n$&$|H_n|$&$|H_n'|$&pairs&gauge orbits&descend&cycles&generating cycles\\
\midrule
2&1&1&1&1&1&1&1\\
3&12&4&8&4&6&8&6\\
4&4&1&4&2&4&2&2\\
5&60&60&20&4&20&20&20\\
6&12&4&8&4&6&8&6\\
7&168&168&42&6&42&42&42\\
8&32&2&16&4&16&4&4\\
12&96&8&24&4&24&12&12\\
24&768&16&96&8&96&24&24\\
\bottomrule
\end{tabular}\end{center}
\begin{proof}
The algorithm constructs $Q_n$ by enumerating determinant-one matrices,
identifying a matrix with its negative, and closing $L_0,L_1$ under
multiplication. It constructs $H_n'$ from all conjugates of $[L_1,L_0]$;
in a two-generator group their normal closure is the derived group.
For each allowed $\nu$ and each $\xi\in H_n$, $\eta$ is forced by
(B.f2.1), after which both products are evaluated exactly. Gauge orbits
use $(\xi,\eta)\mapsto(L_0^a\xi L_1^{-a},L_0^a\eta L_1^{-a})$;
the program checks that every such orbit stays in the solution set.

``Descend'' is a separate test: the proposed images
$S\mapsto\eta S$, $R\mapsto\xi R$ are propagated along the Cayley
graph of $Q_n$ and checked on every edge. Consistency and bijectivity
are both required. Thus it is a test of an automorphism of this
particular finite quotient, not a presumed consequence of two order
relations. The cycle columns enumerate $f\in H_n'$ with the swapping
and cyclic maps actually defined on $H_n$, then test whether
$L_1^{2\nu+1}$ and $f^{-1}L_0^{2\nu+1}f$ generate $H_n$.
All loops are finite and exhaustive; all solutions, not just the counts,
are retained in \nolinkurl{finite_matrix_cycles.json}. Level $1$ is the
single trivial solution. No conclusion about lifting any remaining
residue to $\Fhat$ is built into the algorithm.
\end{proof}
For example, the first failed quotient descent at level $3$ has
$\nu=0$ and $\xi=\eta=S$. It sends $S$ to $1$, and hence is not an
automorphism. This is why Grothendieck's generation condition cannot be
discarded. The outer images among the descending pairs have orders
$1,2,1,2,2,2,4,4,8$ in the displayed level order. The projective level-four
outer image is trivial because $Q_4\simeq S_4$ has no outer automorphisms;
it must not be substituted for the nontrivial outer image on $\SL_2(\mathbb Z/4)$.

\begin{claim}{B.f3}{Proved}{The characteristic $S_3$ calculation}
The characteristic envelope $\overline{S_3}$ has order $108$ and
exponent $6$, with derived group of order $27$. There are exactly twelve
pairs $(\lambda\bmod12,f)$ satisfying (I),(II), with
$f\in\overline{S_3}'$ and an induced automorphism. Their outer image is
\[
 \mathcal D(S_3)=\mathcal D_1(S_3)\simeq C_2.
\]
Complex conjugation is trivial in this particular outer quotient.
\end{claim}
% anchor:B.f3
\manuscriptanchor{\S 42; fols.331 verso--337; cote 140-3 scans85--96}

\begin{proof}
There are three generating-pair orbits modulo $\Aut(S_3)$.
They give an embedding $\overline{S_3}\subseteq S_3^3$.
It is the subgroup in which the product of the three signs is $+1$:
its image on signs has order four, and the squares of the distinguished
generators and their product generate the three independent $C_3$
factors. Thus its order is $4\cdot3^3=108$.

\texttt{finite\_shadows.py} enumerates all four unit residues modulo
$12$ and all $27$ commutator elements. It computes the actual
substitution automorphisms and checks both cycle products, generation,
and the induced homomorphism on every Cayley edge. An independent
enumeration uses the definition of $\mathcal D(S_3)$: fix
$y\mapsto y^k$, range over every conjugate of $x^k$, and check
commutation with both $\theta,\delta$ modulo every possible inner
conjugator. The two enumerations yield the same two outer classes.
The twelve pairs and their orders are printed in the JSON certificate.
The involution $x\mapsto x^{-1},y\mapsto y^{-1}$ has the same outer
key as the identity. This last computation prevents an erroneous
identification of every finite order-two Galois image with complex
conjugation.
\end{proof}

\begin{claim}{B.f4}{Proved}{The arithmetic image of the $S_3$ envelope is all of $C_2$}
The outer Galois image in B.f3 is $C_2$, with fixed field
$\mathbb Q(\sqrt3)$. Consequently the images of $G_\mathbb Q$, $\GT$
and $\GTz$ in $\mathcal D(S_3)$ all coincide; there is no nonarithmetic
outer shadow at this level.
\end{claim}
% anchor:B.f4
\manuscriptanchor{\S 44; fols.365 verso--368 verso; cote 140-3 scans152--158}

\begin{proof}
Use the rational $S_3$-Galois cover
\[
 b(t)=\frac{27t^2(t-1)^2}{4(t^2-t+1)^3},\qquad
 a(t)=\frac{t^3-3t+1}{t(t-1)},\quad w=a-\frac32.
\]
The rational transformations $t\mapsto1-t,1/t$ generate its deck group;
the degree is six. The three distinct orders-of-inertia patterns are
obtained by taking $b=T,1/T,1/(1-T)$, so these are exactly the three
epimorphism kernels used in B.f3. Direct rational identities give
\[
 a\bigl(1/(1-t)\bigr)=a(t),\quad a(1-t)=3-a(t),\quad
 w^2=\frac{27}{4}\frac{1-b}{b}.
\]
Thus the quadratic subfields of the three covers have square classes
$3T(1-T)$, $3(T-1)$, and $-3T$. Their product is the square class of $3$.
Their compositum therefore contains $\sqrt3$ as a constant.
Its geometric Galois group has order $108$ by B.f3, whereas its
arithmetic group is a subgroup of $S_3^3$ of order at most $216$.
The constant quadratic extension forces its order to be at least $216$;
both bounds are equal. No other constant is possible.

The centralizer of the geometric subgroup in $S_3^3$ is trivial,
since its projection to every center-free factor is surjective.
An element outside the geometric subgroup cannot act by an inner
automorphism of it: multiplying it by the inverse of a putative inner
conjugator would centralize that subgroup. Hence the constant-field
quotient acts faithfully in the outer group. This proves the claim.
All rational identities are independently checked in
\texttt{finite\_s3\_fields.py}. The arithmetic character is
$\chi_{-4}\chi_{-3}$, in agreement with C.a2.
\end{proof}

\begin{claim}{B.f5}{Known}{Exact simple-group finite shadows, recomputed}
For the characteristic envelopes attached to $A_5$ and $\PSL_2(\mathbb F_7)$,
\[
 \mathcal D_1(A_5)=1,\quad |\mathcal D(A_5)|=2,
 \qquad
 \mathcal D_1(\PSL_2(\mathbb F_7))\simeq C_2^3\times D_8^2,
 \quad |\mathcal D(\PSL_2(\mathbb F_7))|=1024,
\]
where $D_8$ has order eight. The first of these full shadows is entirely
arithmetic, with cyclotomic quadratic field $\mathbb Q(\sqrt5)$.
We do not determine the order of the arithmetic image in the second.
The full $1024$-element shadow has exactly one element of order one,
$575$ of order two, and $448$ of order four; its exponent is four.
\end{claim}
% anchor:B.f5
\manuscriptanchor{\S 42; fol.349 verso; scan120}

\begin{proof}
We use and verify the mechanism of
\citet[Theorem 4.9 and Proposition 4.12]{Guillot2015root}.
For a nonabelian simple $G$, the distinct epimorphism kernels give
$\overline G=G^r$. This follows by induction from the subdirect-product
lemma: a proper subdirect subgroup of a product with one new simple
factor identifies that factor with an old quotient, contradicting
distinctness of the kernels. Every automorphism of $G^r$ permutes its
simple normal factors and acts on each by $\Aut(G)$.
Consequently its outer group acts faithfully on generating pairs modulo
simultaneous inner conjugacy.

On this set, take the centralizer of the permutations induced by
$\Out(G)$, swapping the pair, and $(x,y)\mapsto((xy)^{-1},y)$.
Intersect it with the permutations preserving the ordered pair of
conjugacy classes. The result is exactly $\mathcal D_1(G)$: the
class-preserving condition translates coordinatewise into conjugacy of
the distinguished elements in $G^r$. The program constructs the set,
these permutations, and every color-preserving equivariant bijection
between their orbits. If a block contains $r$ isomorphic colored orbits
with $a$ equivariant self-maps, its factor has order $a^r r!$.
It also generates every element of the resulting group and checks the
centralizer and color conditions anew.

For $A_5$ the set has $38$ elements (from $2280$ generating pairs), and
the centralizer intersection is trivial. For $\PSL_2(\mathbb F_7)$ it
has $114$ elements (from $19152$ generating pairs); the blocks give
three $C_2$ factors and two $C_2\wr S_2=D_8$ factors. The orders of its
elements are $1$ once, $2$ in $287$ cases, and $4$ in $224$ cases.
These are independent recomputations of the corresponding entries of
\citet[Theorem 5.2]{Guillot2015root}.

The full shadow is obtained by allowing a common power $k$ on all
conjugacy classes. Every unit power pattern is realized by Galois,
because the cyclotomic character is surjective and inertia is raised
to that power; this is also \citet[Lemma 2.3]{Guillot2015root}.
Two such patterns give the same coset precisely when every class is
preserved by their ratio. The power-class kernel has index two in both
examples, as verified for every group element. For $A_5$ it is the
quadratic character modulo five. Since the residual class-preserving
group is trivial, the arithmetic image is the entire order-two group.
For $\PSL_2(\mathbb F_7)$ this argument only gives the quotient of
order two, not all $512$ elements in its kernel.

For the full finite-order calculation, adjoin to the recorded subgroup
$\mathcal D_1$ the permutation
\[
 c:[(x,y)]\longmapsto[(x^{-1},y^{-1})]
\]
of the $114$ inner generating-pair classes. The power class of $-1$
is nontrivial modulo seven, so $c\notin\mathcal D_1$. The established
index-two result identifies the generated group with the full shadow.
The program \nolinkurl{finite_shadow_orders.py} verifies that $c$
normalizes all nine recorded generators, enumerates the full group,
checks all $1024^2$ products for closure, and computes every element
order. The coset $c\mathcal D_1$ contains $288$ involutions and $224$
elements of order four. The other three named full shadows are $C_2$,
so their only orders are one and two. These are finite-shadow orders;
the full arithmetic subgroup at $\PSL_2(\mathbb F_7)$ remains undetermined.
\end{proof}

\subsection{A complete additional finite case: the characteristic $\mathfrak S_4$ envelope}

\begin{claim}{B.s4.1}{Proved}{The characteristic $\mathfrak S_4$ shadow is arithmetic of order two}
Let $N_G$ be the intersection of the kernels of all continuous epimorphisms
$\widehat F_2\twoheadrightarrow G$, and put $\overline G=\widehat F_2/N_G$.
For $G=\mathfrak S_4$, let $\mathrm{GT}(G)$ and $\mathrm{GT}_1(G)$ have
Guillot's finite-shadow definitions (power-conjugacy and commutation with the
two standard outer symmetries, not an asserted image of the full profinite
pentagon group). Then
\[
 |\overline{\mathfrak S_4}|=442368=108\cdot2^{12},\qquad
 \mathrm{GT}(\mathfrak S_4)=\mathrm{GT}_1(\mathfrak S_4)
       =\operatorname{im}(G_{\mathbb Q}\to
           \operatorname{Out}(\overline{\mathfrak S_4}))\simeq C_2.
 \tag{B.s4.1.1}
\]
There is therefore no nonarithmetic finite-shadow element in this quotient.
This supplements the requested finite range of B(e); it asserts no
surjectivity theorem for arbitrary finite shadows.
\end{claim}
% anchor:B.s4.1
\manuscriptanchor{\S 42; fol.349 verso; cote140-3 scan120; added finite-quotient task B(e)}

\begin{proof}
\emph{The envelope and its kernel.}
All automorphisms of $\mathfrak S_4$ are inner. Its $216$ ordered generating
pairs form nine simultaneous-conjugacy classes, represented by the pairs
printed in \texttt{arithmetic\_s4.json}. Write $x,y\in\mathfrak S_4^9$
for their two coordinate tuples and $H=\langle x,y\rangle$. By the
intersection-of-kernels definition, $H=\overline{\mathfrak S_4}$.
Project each factor through $\mathfrak S_4\twoheadrightarrow\mathfrak S_3$
by its action on the three partitions into two pairs. Its kernel is the
Klein four group $V$. The projected group $Q$ has order $108$: each of the
three generating-pair classes of $\mathfrak S_3$ occurs three times among
the nine projections, and after conjugating coordinates it is the
threefold repeated embedding of
\[
 Q_0=\{(a,b,c)\in\mathfrak S_3^3:
           \operatorname{sgn}(a)\operatorname{sgn}(b)
             \operatorname{sgn}(c)=1\}.
 \tag{B.s4.1.2}
\]
The program verifies this by a $108$-element traversal, independent of a
large permutation-group package. A Schreier transversal in $H$ over $Q$
gives $216$ generators for the kernel $V_0=H\cap V^9$. Identify $V^9$
with $\mathbb F_2^{18}$. Their span has rank $12$, with the full bit-vector
basis printed in the output; Schreier's lemma proves that it is exactly
the kernel. Hence $|H|=108\cdot2^{12}$.

\emph{The permutation action has trivial image.}
Let $\mathcal P_c$ be the nine generating-pair classes. The two symmetries
are
\[
 [a,b]\longmapsto[b,a],\qquad
 [a,b]\longmapsto[(ab)^{-1},b].
\]
Their permutations on the program's zero-based indices are respectively
\[
 (2,5,0,6,7,1,3,4,8),\qquad(7,4,5,8,1,2,6,0,3).
\]
The blocks obtained by fixing the separate conjugacy classes of the two
entries are
\[
 \{0\},\{1\},\{2\},\{3,4\},\{5\},\{6,7\},\{8\}.
\]
Of the four block-preserving permutations, only the identity commutes with
both displayed permutations, as direct substitution verifies. Thus
$\mathcal S(\mathfrak S_4)=1$ in the terminology of
\citet[Proposition 4.3, Lemmas 4.5--4.7, \S4.3]{Guillot2015root}.
Every unit power of a permutation has the same cycle type. Consequently
the action of the entire finite $\mathrm{GT}(\mathfrak S_4)$, not just
its multiplier-one subgroup, preserves these blocks. It has trivial image
on $\mathcal P_c$.

\emph{Bounding the kernel of the action.}
Put $N=N_{\mathfrak S_4^9}(H)$ and $W=V^9$.
An automorphism of $H$ acting trivially on $\mathcal P_c$ is, in every
coordinate, conjugation by some element of $\mathfrak S_4$. Conversely
such a tuple defines an automorphism of $H$ exactly when it belongs to
$N$. The centralizer of $H$ in $\mathfrak S_4^9$ is trivial, because its
coordinate projections are surjective and $Z(\mathfrak S_4)=1$.
Modulo inner automorphisms of $H$, the kernel is therefore exactly $N/H$.
In particular $\mathrm{GT}(\mathfrak S_4)\subseteq N/H$.

We compute this normalizer quotient without enumerating its ambient group.
The normalizer of the repeated group $Q$ in $\mathfrak S_3^9$ lies in
the corresponding repeated copy of $\mathfrak S_3^3$: within any repeated
block, two conjugators must induce the same conjugation on the surjective
$\mathfrak S_3$ factor, and so coincide, since its centre is trivial.
By (B.s4.1.2) the resulting quotient of normalizer by $Q$ has order two.
Next
\[
 (W/V_0)^Q=0. \tag{B.s4.1.3}
\]
There is a conceptual check as well as the rank certificate: the normal
odd-order subgroup $A=C_3^3$ of $Q_0$ projects onto $C_3$ in every
coordinate, which acts without fixed vectors on each copy of
$V\simeq\mathbb F_2^2$. Hence $W^A=0$. Taking invariants under $A$ is
exact in characteristic two, by averaging over its odd order, and so
$(W/V_0)^A=0$ and a fortiori (B.s4.1.3). The program independently
computes the two-generator fixed-vector equations: their rank is $6$ on
$W$, their solution space has dimension $12$, and quotienting by the
$12$-dimensional $V_0$ leaves dimension zero.
If $w\in N\cap W$, then $[w,H]\subseteq V_0$, so (B.s4.1.3) implies
$w\in V_0$. Thus $N\cap W=V_0$ and
\[
 N/H\hookrightarrow N_{\mathfrak S_3^9}(Q)/Q\simeq C_2.
\]

\emph{An explicit nontrivial finite-shadow class.}
The code solves the normalizer equations in the $64$-element vector-space
quotient $W/V_0$, and prints a tuple $n\in N\setminus H$. In the nine
coordinates and using permutation arrays on $\{0,1,2,3\}$, it is
\[
 n=(1,(01),1,1,1,(13),1,1,(13)).
 \tag{B.s4.1.4}
\]
Here $1$ denotes the identity permutation; the cycle notation is in the
zero-based labeling just specified. The output additionally prints
conjugators \emph{in $H$} showing that $n^{-1}xn$ is conjugate to $x$
and $n^{-1}yn$ to $y$. These tests use the Schreier transversal and exact
linear solving in $V_0$, and verify the resulting conjugators by direct
multiplication. The kernel $N/H$, being the kernel of the natural action
of $\operatorname{Out}(H)$ on $\mathcal P_c$, is normal in that outer
automorphism group. Since it has order two, all outer automorphisms
centralize it, in particular the two standard symmetries. Therefore the
class of $n$ is in $\mathrm{GT}_1(\mathfrak S_4)$, proving both finite
shadow orders in (B.s4.1.1). We emphasize the extra conjugacy check:
rational conjugacy classes of $\mathfrak S_4$ alone would not prove
conjugacy by an element of the proper subgroup $H\subset\mathfrak S_4^9$.

\emph{The arithmetic lower bound.}
The projection to $Q$ identifies it with
$\overline{\mathfrak S_3}$, so $N_{\mathfrak S_4}\subseteq
N_{\mathfrak S_3}$ and the arithmetic outer action descends. We verify the
nontriviality of this descended action by the explicit fields, rather than
assuming every finite shadow is arithmetic. The rational quotient map
\[
 \beta(t)=\frac{27t^2(t-1)^2}{4(t^2-t+1)^3}
\]
is invariant under $t\mapsto1-t$ and $t\mapsto1/t$, has degree six,
and is a regular $\mathfrak S_3$ cover. Its branch orders at $0,1,\infty$
are $(2,2,3)$. The covers obtained by the equations
$\beta(t)=T$, $\beta(t)=1/T$ and $\beta(t)=1/(1-T)$ therefore represent
all three generating-pair classes. If
\[
 w=\frac{t^3-3t+1}{t(t-1)}-\frac32,
\]
then $w^2=\frac{27}{4}(1-\beta)/\beta$; this rational-function identity
and the invariances are verified in \texttt{finite\_s3\_fields.py}.
The three quadratic subfields have square classes
\[
 3T(1-T),\quad3(T-1),\quad-3T.
\]
Their product is $27T^2(T-1)^2$, so their compositum contains
$\mathbb Q(\sqrt3)$. Its geometric group has the already computed order
$108$, while its arithmetic group is a subgroup of $\mathfrak S_3^3$,
of order at most $216$. The degree of its constant field is consequently
at most two and at least two; it is exactly $\mathbb Q(\sqrt3)$, and
the arithmetic group is all of $\mathfrak S_3^3$. An element outside the
geometric subgroup acts noninner on it: otherwise the element divided
by an appropriate geometric conjugator would centralize that subgroup,
whose centralizer in $\mathfrak S_3^3$ is trivial by the surjective
coordinate projections. Thus the arithmetic outer image on
$\overline{\mathfrak S_3}$ has order two. The image on $H$ has order
at least two, and is contained in the finite shadow of order two just
proved. This finishes (B.s4.1.1).
\end{proof}

\begin{claim}{B.s4.2}{Proved}{What the simple-group shadow counts do and do not prove}
The finite calculations for $A_5$ and $\mathrm{PSL}_2(\mathbb F_7)$,
when interpreted by \citet[Theorem 4.9]{Guillot2015root}, give
\[
\begin{array}{c|rrr}
 G&|\mathrm{GT}_1(G)|&|\mathrm{GT}(G)|&
       |\operatorname{im}G_{\mathbb Q}|\\\hline
 A_5&1&2&2\\
 \mathrm{PSL}_2(\mathbb F_7)&512&1024&\text{not determined here}
\end{array}
\]
The lower bound on the last unknown size is two. Neither the colored
centralizer computation nor cyclotomic surjectivity proves its equality
to $1024$.
\end{claim}
% anchor:B.s4.2
\manuscriptanchor{\S 42; fol.349 verso; cote140-3 scan120; added finite-quotient task B(e)}

\begin{proof}
Both groups are finite, nonabelian and simple, exactly the hypothesis of
Guillot's Theorem 4.9. The standard outer involution is conjugation by an
odd permutation for $A_5$, and by a matrix of nonsquare determinant for
$\mathrm{PSL}_2(\mathbb F_7)$. These generate their respective outer
automorphism groups of order two; this is also verified by enumerating
all automorphisms in \texttt{arithmetic\_simple\_audit.py}. The
permutations used in \texttt{finite\_shadows.py} are exactly the swap,
$(g,h)\mapsto((gh)^{-1},h)$, and those outer involutions appearing in
the theorem. The color is the pair of separate conjugacy classes, not
an incorrectly assumed action of $G\times G$ on generating pairs.
Thus the computed colored centralizer orders $1$ and $512$ have precisely
the stated multiplier-one meaning.

For completeness we justify the passage to the full finite group.
Guillot's Lemma 4.10 gives $\overline G=G^r$. Therefore $x^k$ is
conjugate to $x$ in $\overline G$ exactly when each coordinate has the
same conjugacy class under the power $k$. In these two groups every
nontrivial conjugacy class occurs in a generating pair; this too is
checked in the audit. The kernel of the power action on all conjugacy
classes is
\[
\begin{array}{c|l}
 A_5,\ e=30&1,11,19,29\pmod {30}\\
 \mathrm{PSL}_2(\mathbb F_7),\ e=84&
 1,11,23,25,29,37,43,53,65,67,71,79\pmod {84}.
\end{array}
\]
Its index among all units is two in each case. The quotient of the full
finite shadow by its multiplier-one subgroup embeds in this power-action
group, by its definition, and the cyclotomic character of Galois realizes
every residue unit. Hence it is all of that group and doubles the orders.
For $A_5$ this same arithmetic lower bound fills the entire group of
order two. For $\mathrm{PSL}_2(\mathbb F_7)$ it says nothing about how
much of the kernel of order $512$ is realized by Galois. The smallest
unsettled question within this row remains the exact image on this
specified characteristic quotient, not the polynomial-time computation
of its finite-shadow upper bound.
\end{proof}


\begin{claim}{B.f6}{Proved}{A noncongruence finite quotient detects information lost by matrices}
The continuous map
$\widehat{\SL_2(\mathbb Z)}\to\SL_2(\Zhat)$ is not injective.
The corresponding loss is already present in the index-six free
subgroup after passing to projective groups. Thus the matrix equations
are necessary conditions and are not an equivalent encoding of the
full profinite equations merely because their matrix calculation is correct.
\end{claim}
% anchor:B.f6
\manuscriptanchor{\S \S 39,42; fols.299--301,349 verso; scans22--26,120}

\begin{proof}
Let $\mathbb F_8=\mathbb F_2[a]/(a^3+a+1)$. The matrices
\[
 A=\begin{pmatrix}1&1\\0&1\end{pmatrix},\qquad
 B=\begin{pmatrix}0&a\\a^{-1}&1\end{pmatrix}
\]
have orders two and three and generate $\SL_2(\mathbb F_8)$ of order
$504$. The program \nolinkurl{finite_noncongruence.py} verifies generation
and simplicity: for every nonidentity conjugacy class its normal
closure is the whole group. This is a complete finite simplicity
certificate, not a probabilistic test. Hence $C_2*C_3$ surjects onto
this group, and the map extends to its profinite completion.

No congruence quotient has this simple quotient. By the Chinese
remainder theorem and the successive reduction kernels, the
nonabelian simple composition factors of $\SL_2(\mathbb Z/n)$ are
$\PSL_2(\mathbb F_p)$ for primes $p\ge5$; the kernels between prime
powers are $p$-groups and the cases $p=2,3$ are solvable. The order
$p(p^2-1)/2$ is $168$ at $p=7$ and $660$ at $p=11$, strictly increasing,
and never $504$. A continuous map from the congruence completion to a
finite group factors through some finite level. Thus the quotient of
order $504$ cannot factor through that completion. If the completion
map were injective it would be an isomorphism of compact groups, which
would force just such a factorization.

Its projective kernel lies in the kernel of reduction modulo two,
which is the free subgroup $F$. If the comparison restricted to $F$
were injective, that projective kernel would vanish. This proves the
free-group version as well. What is lost is the closed congruence
kernel and the action on all its finite quotients, not merely a sign.
\end{proof}

\begin{claim}{B.f7}{Open}{The exact remaining finite-image and pentagon problem}
For $G=S_3,S_4,A_5$, the outer characteristic shadows of $G_\mathbb Q$,
$\GT$, and $\GTz$ are all the computed order-two group. For
$G=\PSL_2(\mathbb F_7)$, the unresolved arithmetic image is a subgroup
of the explicitly given group of order $1024$, surjecting onto its
power-class quotient $C_2$. Determining this subgroup requires the
Galois action on the $57$ regular covers and on their labeled structures,
not the conjugacy-class calculation alone. In addition, the raw
matrix residues in B.f2 do not come with the compatible five-puncture
quotient required to test (III); the uncomputed datum is that quotient
and its five evaluation maps, already at level $n=3$, the first nontrivial Legendre matrix image
after the trivial level-two image.
\end{claim}
% anchor:B.f7
\manuscriptanchor{\S 42; fol.349 verso; scan120}

The attack is B.f2--B.f5 and B.s4.1: complete residue lists and finite
shadow groups are available, and arithmetic fields close three of the
four named cases. No finite shadow outside the arithmetic image has
been established in this range. We have neither equated a finite shadow
with the image of the inverse limit nor reported a pentagon count for
an undefined finite target. The first unresolved member of the four
named characteristic examples is the group of order $168$; a universal
characteristic-envelope pentagon calculation for all the separately
listed $\SL_2(\mathbb Z/n)$ is not supplied by these matrix tests.

\begin{claim}{B.f8}{Proved}{The BISITE endomorphism is invisible to every nilpotent completion}
The endomorphism $v:F_2\to F_2$ defined by
$v(x)=x$, $v(y)=y[x,y]$ induces an automorphism of every lower-central
nilpotent quotient, every free pro-$\ell$ completion, and every
characteristic-zero pro-unipotent completion. It is not an automorphism
of the full profinite completion. It preserves the first two peripheral
conjugacy classes, but not the third; thus it is not a counterexample
inside $\GTz$ or the group of all loop automorphisms.
\end{claim}
% anchor:B.f8
\manuscriptanchor{\S 42; fol.349 verso; cote140-3 scan120; added BISITE B.11 follow-up}

\begin{proof}
It is the identity on abelianization. The graded Lie ring of a free
group is generated by degree one, so the induced maps on every
lower-central graded quotient are the identity. Induction on the
central extensions proves the nilpotent quotient assertion. For the
pro-$\ell$ completion the images generate the Frattini quotient,
so they topologically generate the free pro-$\ell$ group. The resulting
surjection is an automorphism: a finitely generated profinite group is
Hopfian, since it has only finitely many open normal subgroups of any
fixed index, and a surjective endomorphism permutes these subgroups.
Alternatively the Magnus realization and the preceding nilpotent
argument prove the same result. In the pro-unipotent Lie algebra the
map is a filtered identity on the associated graded, and is inverted
recursively degree by degree.

On three letters put $x=(12)$, $y=(13)$. Exact multiplication gives
$y[x,y]=(12)$. The old pair generates $S_3$ and is transitive, while
the new pair generates $C_2$ and fixes the third letter. Precomposing
a transitive finite action with a surjective profinite endomorphism
would preserve its image and its orbits. Hence the profinite map is
not surjective. The full calculation is in
\texttt{finite\_congruence.json}.
Finally $y[x,y]=(yx)y(yx)^{-1}$ is conjugate to $y$, but
$v(xy)=xyxyx^{-1}y^{-1}$ is cyclically reduced of length six, whereas
$xy$ has cyclic length two. They are not conjugate in the discrete
free group. In the same finite $S_3$ quotient, $xy$ has order three
whereas its image under $v$ is the identity, so they are not conjugate
in the full profinite group either. The third-loop condition cannot therefore be inferred
from the first two, or from the completion assertions just proved.
\end{proof}
